% Source: https://artofproblemsolving.com/community/c3h2264405_mathcounts_year_2017_sprint__23
% Stable identity: aops-topic-2264405-q01
Write the number as $N=\overline{AB28}$ with digits $A\ne0$ and $B$. Its reversal is $R=\overline{82BA}=8200+\overline{BA}$.

\textbf{The reversed number.} $8200=16\cdot512+8$, so $R=16\cdot512+\left(8+\overline{BA}\right)$, and $R$ is a multiple of $16$ exactly when $8+\overline{BA}$ is, that is, when $\overline{BA}$ leaves remainder $16-8=8$ on division by $16$. The two-digit numbers (with leading zero allowed) of this form are $08,24,40,56,72,88$. Since $A$ is the leading digit of $N$, $A\ne0$, which rules out $40$.

\textbf{The number itself.} $N=100\cdot\overline{AB}+28$. Since $100=16\cdot6+4$ and $28=16+12$, $N$ leaves the same remainder as $4\cdot\overline{AB}+12$, so $16\mid N$ exactly when $4\cdot\overline{AB}\equiv4\pmod{16}$, that is, $\overline{AB}\equiv1\pmod4$.

Check the five candidates $(B,A)$: $\overline{AB}=80,42,65,27,88$. Only $65\equiv1\pmod4$. So $A=6$, $B=5$, and the number is $\boxed{6528}$. Indeed $6528=16\cdot408$ and $8256=16\cdot516$.
